\documentclass[10pt]{article}

\usepackage{amsmath,amsthm}

\newtheorem{definition}{Definition}

\title{Faster Triangle Detection in CONGEST -- Hard}
\author{Keren Censor-Hillel}
\date{}

\begin{document}

\maketitle

\section{Preliminaries}

\begin{definition}[CONGEST model]
An input is a simple undirected graph $G=(V,E)$ with $|V|=n$, which is also
the communication network. The vertices have distinct $O(\!\log n)$-bit
identifiers, and computation proceeds in synchronous rounds. In each round,
the endpoints of every edge can exchange an $O(\!\log n)$-bit message in
each direction. Initially, each vertex knows $n$, its own identifier, and
the identifiers of its neighbors in $G$. Local computation is unrestricted.
\end{definition}

\begin{definition}[Triangle detection]
A triangle in $G$ consists of three distinct vertices $u,v,w$ such that
$uv,vw,wu\in E$. An algorithm detects triangles if all vertices output $0$
whenever $G$ is triangle-free, while at least one vertex outputs $1$
whenever $G$ contains a triangle.

A randomized $T(n)$-round algorithm must terminate within $T(n)$ rounds on
every execution and satisfy the preceding condition with probability at
least $2/3$ on every input graph and every assignment of identifiers.
\end{definition}

The benchmark upper bound for general graphs is
$\widetilde O(n^{1/3})$ rounds, where $\widetilde O$ suppresses
polylogarithmic factors in $n$.

\section{Problem formulation}

Prove that there exist constants $0<\varepsilon\leq 1/3$ and $C\geq 0$,
and a possibly randomized CONGEST algorithm with error at most $1/3$, that
detects triangles in
\[
O\!\left(n^{1/3-\varepsilon}(\log n)^C\right)
\]
rounds.

\end{document}
